\subsection{BÀI TẬP TRẮC NGHIỆM}
\ind{PHẦN I.} \inden{Câu trắc nghiệm nhiều phương án lựa chọn. Mỗi câu hỏi học sinh chỉ chọn một phương án.}\\
\setcounter{ex}{0}
\Opensolutionfile{ans}[ans/1H7-B4-1]
\begin{ex}%[1H3B4-1]
	Trong khẳng định sau về lăng trụ đều, khẳng định nào \textbf{sai}?
	\choice
	{Đáy là đa giác đều}
	{Các mặt bên là những hình chữ nhật nằm trong mặt phẳng vuông góc với đáy}
	{\True Các mặt bên là những hình vuông}
	{Các cạnh bên là những đường cao}
	\loigiai{
		Vì lăng trụ đều là lăng trụ đứng nên các cạnh bên bằng nhau và cùng vuông góc với đáy. Do đó các mặt bên là những hình chữ nhật.
	}
\end{ex}

\begin{ex}%[1H3B4-1]
	Trong các mệnh đề sau, mệnh đề nào đúng?
	\choice
	{Nếu hình hộp có bốn đường chéo bằng nhau thì nó là hình lập phương}
	{Nếu hình hộp có sau mặt bằng nhau thì nó là hình lập phương}
	{Nếu hình hộp có hai mặt là hình vuông thì nó là hình lập phương}
	{\True Nếu hình hộp có ba mặt chung một đỉnh là hình vuông thì nó là hình lập phương}
	\loigiai{
		Nếu hình hộp có ba mặt chung một đỉnh là hình vuông thì nó là hình lập phương.
	}
\end{ex}

\begin{ex}%[1H3Y4-1]
	Khẳng định nào sau đây đúng?
	\choice
	{\True Hình lăng trụ đứng có đáy là một đa giác đều là hình lăng trụ đều}
	{Hình lăng trụ đứng là hình lăng trụ đều}
	{Hình lăng trụ tứ giác đều là hình lập phương}
	{Hình lăng trụ có đáy là một đa giác đều là hình lăng trụ đều}
	\loigiai{
		Hình lăng trụ đứng có đáy là một đa giác đều là hình lăng trụ đều.
	}
\end{ex}

\begin{ex}
	Khẳng định nào sau đây \textbf{sai} khi nói về hình chóp tứ giác đều?
	\choice
	{Các cạnh bên bằng nhau}
	{\True Các mặt bên là những tam giác đều}
	{Các cạnh đáy bằng nhau }
	{ Mặt đáy là hình vuông}
\end{ex}

\begin{ex}%[1H3Y4-2]
	Cho hình lăng trụ đứng $ABC.A'B'C'$ có đáy $ABC$ là tam giác vuông cân tại $A$, chiều cao của lăng trụ là cạnh
	\choice
	{\True $BB'$}
	{$AC'$}
	{ $AB$}
	{$AB'$}
	\loigiai{
		Theo định nghĩa về lăng trụ đứng là lăng trụ có cạnh bên vuông góc với đáy nên ta chọn đáp án là cạnh $BB'$ là đường cao của lăng trụ đứng $ABC.A'B'C'$.}
\end{ex}

\begin{ex}%[1H3B4-1]
	Trong các mệnh đề sau, mệnh đề nào sau đây là đúng?
	\choice
	{\True Hai mặt phẳng vuông góc với nhau thì mọi đường thẳng nằm trong mặt phẳng này và vuông góc với giao tuyến của hai mặt phẳng sẽ vuông góc với mặt phẳng kia}
	{Hai mặt phẳng vuông góc với nhau thì mọi đường thẳng nằm trong mặt phẳng này sẽ vuông góc với mặt phẳng kia}
	{Hai mặt phẳng phân biệt cùng vuông góc với một mặt phẳng thì vuông góc với nhau}
	{Hai mặt phẳng phân biệt cùng vuông góc với một mặt phẳng thì song song với nhau}
	\loigiai{
		Mệnh đề ``Hai mặt phẳng vuông góc với nhau thì mọi đường thẳng nằm trong mặt phẳng này sẽ vuông góc với mặt phẳng kia'' là sai. Hai mặt phẳng vuông góc với nhau thì đường thẳng nằm trong mặt phẳng này, vuông góc với giao tuyến thì vuông góc với mặt phẳng kia.\\
		Mệnh đề ``Hai mặt phẳng phân biệt cùng vuông góc với một mặt phẳng thì vuông góc với nhau'' và ``Hai mặt phẳng phân biệt cùng vuông góc với một mặt phẳng thì song song với nhau'' là sai. Hai mặt phẳng phân biệt cùng vuông góc với một mặt phẳng thì song song với nhau hoặc cắt nhau (giao truyến vuông góc với mặt phẳng kia).
	}
\end{ex}

\begin{ex}%[1H3B4-1]
	Trong các mệnh đề sau, mệnh đề nào đúng?
	\choice
	{Qua một đường thẳng có duy nhất một mặt phẳng vuông góc với một đường thẳng cho trước}
	{Qua một điểm có duy nhất một mặt phẳng vuông góc với một mặt phẳng cho trước}
	{Hai mặt phẳng phân biệt cùng vuông góc với một mặt phẳng thì song song với nhau}
	{\True Hai mặt phẳng phân biệt cùng vuông góc với một đường thẳng thì song song với nhau}
	\loigiai{
		Mệnh đề ``Hai mặt phẳng phân biệt cùng vuông góc với một mặt phẳng thì song song với nhau'' là sai. Hai mặt phẳng phân biệt cùng vuông góc với một mặt phẳng thì song song với nhau hoặc cắt nhau (giao tuyến vuông góc với mặt phẳng thứ 3).\\
		Mệnh đề ``Qua một đường thẳng có duy nhất một mặt phẳng vuông góc với một đường thẳng cho trước'' là sai. Qua một đường thẳng vô số mặt phẳng vuông góc với một đường thẳng cho trước.\\
		Mệnh đề ``Qua một điểm có duy nhất một mặt phẳng vuông góc với một mặt phẳng cho trước'' là sai. Qua một điểm có vô số mặt phẳng vuông góc với một mặt phẳng cho trước.
	}
\end{ex}

\begin{ex}%[TEX SBT CTST, Đình Nguyên]%[1H2Y4-1]
	Hai mặt phẳng phân biệt cùng vuông góc với mặt phẳng thứ ba thì
	\choice
	{Song song với nhau}
	{Trùng nhau}
	{Không song song với nhau}
	{\True Hoặc song song với nhau hoặc cắt nhau theo giao tuyến vuông góc với mặt phẳng thứ ba}
	\loigiai{
		Hai mặt phẳng phân biệt cùng vuông góc với mặt phẳng thứ ba thì hoặc song song với nhau hoặc cắt nhau theo giao tuyến vuông góc với mặt phẳng thứ ba.	
	}
\end{ex}

\begin{ex}%[1H3Y4-3]
	Cho hình lập phương $ABCD.A'B'C'D'$. Góc giữa hai mặt phẳng $(ABCD)$ và $(A'B'C'D')$ bằng bao nhiêu?
	\choice
	{$60^{\circ}$}
	{\True $0^{\circ}$}
	{$45^{\circ}$}
	{$90^{\circ}$}
	\loigiai{
		Hai mặt phẳng $(ABCD)$ và $(A'B'C'D')$ song song nên góc giữa chúng bằng $0^{\circ}$.
	}
\end{ex}

\begin{ex}%[1H3Y4-3]%
	Cho hình chóp $S.ABC$ có $SA \perp (ABC)$ và $AB \perp BC$, gọi $I$ là trung điểm $BC$. Góc giữa hai mặt phẳng $(SBC)$ và $(ABC)$ là góc nào sau đây?
	\choice
	{$\widehat{SCA}$}
	{\True $\widehat{SBA}$}
	{$\widehat{SIA}$}
	{$\widehat{SCB}$}
	\loigiai{
		\immini {Ta có: $BC \perp SA,BC \perp AB \Rightarrow BC \perp SB$\\
			$\Rightarrow \heva{&(SBC) \cap (ABC)=BC\\&AB \perp BC,AB \subset (ABC)\\&SB \perp BC,SB \subset (SBC)}$ $\Rightarrow \widehat{((SBC),(ABC))}=\widehat{SBA}$.}
		{\begin{tikzpicture}[line join=round,line cap=round,line width=.6pt,font=\footnotesize,scale=0.6]
				\coordinate[label=left:$A$] (A) at (0,0);
				\coordinate[label=below left:$B$] (B) at (1.2,-1.8);
				\coordinate[label=right:$C$] (C) at (5,0);
				\coordinate[label=above left:$S$] (S) at ($(A)+(90:4)$);
				\coordinate[label=above right:$I$] (I) at ($(B)!.5!(C)$);
				\draw ($ (A)!7pt!(C)$)--($ ( $(A)!7pt!(C)$ )+ ( $(A)!7pt!(S)$ )-(A) $)--( $(A)!7pt!(S)$);
				\draw ($ (A)!7pt!(B)$)--($ ( $(A)!7pt!(B)$ )+ ( $(A)!7pt!(S)$ )-(A) $)--( $(A)!7pt!(S)$);
				\draw ($ (B)!7pt!(A)$)--($ ( $(B)!7pt!(A)$ )+ ( $(B)!7pt!(C)$ )-(B) $)--( $(B)!7pt!(C)$);
				\draw (A)--(B)--(C)--(S)--cycle (S)--(B) (S)--(I) ;
				\draw[dashed] (A)--(C) (A)--(I) ;
				\foreach \x in {A,B,C,S} \fill[black] (\x) circle (1.5pt);
	\end{tikzpicture}}}
\end{ex}

\begin{ex}%[1H3B4-2]%
	Cho hình chóp $S.ABC$ có đáy $ABC$ là tam giác vuông cân tại $B$, $SA$ vuông góc với đáy. Gọi $M$ là trung điểm của $AC$. Khẳng định nào sau đây \textbf{sai}?
	\choice
	{\True $(SBC)\perp (SAC)$}
	{$(SBM)\perp (SAC)$}
	{$BM\perp AC$}
	{$(SAB)\perp (SBC)$}
	\loigiai{
		\immini{
			\begin{itemize}
				\item Ta có $BC\perp AB$ và $BC\perp SA$ nên $BC\perp (SAB)$, suy ra $(SAB)\perp (SBC)$.
				\item Do $ABC$ là tam giác cân tại $B$ nên $BM\perp AC$.
				\item Ta có $BM\perp SA$, $BM\perp AC$ nên $BM\perp (SAC)$, suy ra $(SBM)\perp (SAC)$.
			\end{itemize}
		}
		{
			\begin{tikzpicture}[>=stealth]
				\tkzDefPoints{0/0/A, 0/4/S, 2/-1/B, 5/0/C}
				\tkzMarkRightAngle(A,B,C)
				\tkzDefMidPoint(A,C)\tkzGetPoint{M}
				\tkzDrawSegments[](S,A S,B S,C A,B B,C)
				\tkzDrawSegments[dashed](A,C B,M S,M)
				\tkzLabelPoints[above](S)
				\tkzLabelPoints[left](A)
				\tkzLabelPoints[below](B)
				\tkzLabelPoints[right](C)
				\tkzLabelPoints[above right](M)
				\tkzDrawPoints[fill=black](S,A,B,C,M)
			\end{tikzpicture}
		}
	}
\end{ex}

\begin{ex}%[1H3B4-2]
	Cho hình chóp tứ giác đều $ S.ABCD $ có $ O $ là giao điểm của $ AC  $ và $ BD $, cạnh bên và cạnh đáy đều bằng $ a $. Khẳng định nào sau đây đúng?
	\choice
	{\True $( SAC )\perp ( ABCD )$}
	{$( SBC )\perp ( SAC )$}
	{$( SAC )\perp ( SCD )$}
	{$( SAD )\perp ( SBD )$}
	\loigiai{
		\begin{center}
			\begin{tikzpicture}
				\tkzDefPoints{0/0/O, -4/-1/A, 1/-1/B, 0/5/S}
				\tkzDefPointBy[symmetry = center O](A)
				\tkzGetPoint{C}
				\tkzDefPointBy[symmetry = center O](B)
				\tkzGetPoint{D}
				\tkzDrawSegments[dashed](A,D A,C D,C D,B S,D S,O)
				\tkzDrawSegments(A,B B,C S,A S,B S,C)
				\tkzMarkRightAngle[size=.2](B,O,C)
				\tkzLabelPoints[below](O,A,B)
				\tkzLabelPoints[above](S)
				\tkzLabelPoints[right](C)
				\tkzLabelPoints[left](D)
				\tkzMarkSegments[mark=|](S,A S,B S,C S,D A,D D,C C,B B,A)
			\end{tikzpicture}
		\end{center}
		Ta có  $SO\perp ( ABCD ),SO\subset ( SAC )$ nên $( SAC )\perp ( ABCD )$}
\end{ex}


\begin{ex}%[Trần Quốc Tráng]%[1H3Y4-2]%
	Cho hình chóp $S.ABCD$ có đáy $ABCD$ là hình vuông, hình chiếu vuông góc của $S$ trên mặt phẳng $(ABCD)$ trùng với trung điểm cạnh $AB$. Khẳng định nào sau đây là \textbf{sai}?
	\choice
	{$(SBC)\perp (SAB)$}
	{$(SAB)\perp (ABCD)$}
	{\True $(SBD)\perp (ABCD)$}
	{$(SHC)\perp (ABCD)$}
	\loigiai{
		\immini{
			Vì $H$ là hình chiếu vuông góc của $S$ trên mặt phẳng $(ABCD)$ nên $SH\perp (ABCD)$.\\
			Các mặt phẳng $(SHC),(SAB)$ đều chứa $SH$ nên chúng vuông góc mặt phẳng đáy $(ABCD)$.\\
			Dễ thấy $BC\perp (SAB)$ nên $(SBC)\perp (SAB)$.
			Do đó các câu $B,C,D$ là đúng.\\
			Suy ra câu $A$ sai.
		}
		{
			\begin{tikzpicture}[scale=.6, font=\footnotesize, line join=round, line cap=round, >=stealth]
				\tkzDefPoints{0/0/B, 7/0/C, 3/3/A, 10/3/D}
				\coordinate (H) at ($(A)!0.5!(B)$);
				\coordinate (S) at ($(H)+(0,6)$);
				\draw (A) node[left]{$A$} circle (1pt);
				\draw (B) node[left]{$B$} circle (1pt);
				\draw (C) node[right]{$C$} circle (1pt);
				\draw (D) node[right]{$D$} circle (1pt);
				\draw (H) node[below right]{$H$} circle (1pt);
				\draw (S) node[above]{$S$} circle (1pt);
				\tkzDrawPolygon(S,C,D)
				\draw (C)--(B)--(S);
				\draw[dashed](S)--(A)--(B)--(D)--(A)--(C) (S)--(H);
			\end{tikzpicture}
		}
	}
\end{ex}

\begin{ex}%[1H2K6-2]
	Cho hình chóp $S. A B C$ có $S A \perp(A B C)$. Gọi $I$ là hình chiếu của $A$ trên $B C$, $\alpha$ là góc giữa đường thẳng $S I$ và mặt phẳng $(A B C)$, $\beta$ là số đo của góc nhị diện $[S, B C, A]$. Phát biểu nào sau đây là đúng?
	\choice
	{$\alpha=90^\circ-\beta$}
	{$\alpha=180^\circ-\beta$}
	{$\alpha=90^\circ+\beta$}
	{\True $\alpha=\beta$}
	\loigiai{
		\immini{
			Ta có $SA\perp (ABC)$ nên
			\[\left(SI,(ABC)\right)=\left(SI,AI\right)=\widehat{SIA}=\alpha. \]
			Lại có $\heva{&BC\perp SA\ (SA\perp (ABC))\\&BC\perp IA\ (\text{giả thiết})}\Rightarrow BC\perp (SAI)$, suy ra $BC\perp SI$. Do đó
			góc nhị diện $[S, B C, A]$ bằng $\widehat{SIA}$.\\
			Vậy $\beta=\widehat{SIA}=\alpha$.
		}{
			\begin{tikzpicture}[>=stealth,line join=round,line cap=round,font=\footnotesize,scale=1]
				\def\a{5};
				\def\b{2};
				\def\g{-40};
				\def\cao{4};
				\path
				(0,0) coordinate (A)
				(0:\a) coordinate (C)
				(\g:\b) coordinate (B)
				($(A)+(90:\cao)$) coordinate (S)
				($(B)!0.2!(C)$) coordinate (I)
				;
				\draw (S)--(A)--(B)--(C)--(S)--(B) (S)--(I);
				\draw[dashed] (I)--(A)--(C);
				\foreach \x/\g in {A/180,B/-90,C/0,S/90,I/-45}\fill[black](\x) circle (1pt) +(\g:3mm) node {$\x$};
				\pic[draw,angle radius=2mm]{right angle=A--I--C};
				\pic[draw,angle radius=6mm,"\tiny $\alpha$"]{angle=S--I--A};
			\end{tikzpicture}
		}
	}
\end{ex}

\begin{ex}%[1H2G6-2]
	Cho hình chóp $S. A B C$ có $S A \perp(A B C)$, $A B \perp B C$, $S A=A B=3 a$, $B C=4 a$. Gọi $\alpha$ là số đo của các góc nhị diện $[A, B C, S]$. Tính $\cos \alpha$.
	\choice
	{$\cos \alpha=\dfrac{1}{2}$}
	{$\cos \alpha=\dfrac{3}{5}$}
	{$\cos \alpha=\dfrac{1}{5}$}
	{\True $\cos \alpha=\dfrac{\sqrt{2}}{2}$}
	\loigiai{
		\immini{
			Ta có $B C \perp(S A B)$ nên $B C \perp S B$ suy ra góc $S B A$ là góc phẳng nhị diện của $[A, B C, S]$. Như vậy, ta có
			\[S B=\sqrt{A B^2+S A^2}=\sqrt{(3 a)^2+(3 a)^2}=3 \sqrt{2} a. \]
			và $\cos \alpha=\cos \widehat{S B A}=\dfrac{A B}{S B}=\dfrac{3 a}{3 \sqrt{2} a}=\dfrac{\sqrt{2}}{2}$.
		}{
			\begin{tikzpicture}[>=stealth,line join=round,line cap=round,font=\footnotesize,scale=1]
				\def\a{5};
				\def\b{2};
				\def\g{-40};
				\def\cao{4};
				\path
				(0,0) coordinate (A)
				(0:\a) coordinate (C)
				(\g:\b) coordinate (B)
				($(A)+(90:\cao)$) coordinate (S)
				($(S)!0.4!(B)$) coordinate (H)
				($(S)!0.6!(C)$) coordinate (K)
				;
				\draw (S)--(A)--(B)--(C)--(S)--(B) (A)--(H)--(K);
				\draw[dashed] (A)--(C) (A)--(K);
				\foreach \x/\g in {A/180,B/-90,C/0,S/90,H/45,K/30}\fill[black](\x) circle (1pt) +(\g:3mm) node {$\x$};
				\pic[draw,angle radius=2mm]{right angle=A--B--C};
				\pic[draw,angle radius=2mm]{right angle=A--H--S};
				\pic[draw,angle radius=2mm]{right angle=A--K--S};
			\end{tikzpicture}
		}
	}
\end{ex}

\begin{ex}%[1H2G6-2]
	Cho hình chóp $S. A B C$ có $S A \perp(A B C)$, $A B \perp B C$, $S A=A B=3 a$, $B C=4 a$. Gọi $\alpha$ là số đo của các góc nhị diện $[B, S A, C]$. Tính $\cos \alpha$.
	\choice
	{$\cos \alpha=\dfrac{1}{2}$}
	{\True $\cos \alpha=\dfrac{3}{5}$}
	{$\cos \alpha=\dfrac{1}{5}$}
	{$\cos \alpha=\dfrac{\sqrt{2}}{5}$}
	\loigiai{
		\immini{
			Vì $S A \perp(A B C)$, $A B \subset(A B C)$, $A C \subset(A B C)$ nên $S A \perp A B$, $S A \perp A C$. Suy ra góc $B A C$ là góc phẳng nhị diện của $[B, S A, C]$, hay $\widehat{B A C}=\alpha$. Xét tam giác $A B C$ vuông tại $B$ có
			\[A C=\sqrt{A B^2+B C^2}=\sqrt{(3 a)^2+(4 a)^2}=5 a. \]
			và $\cos \alpha=\cos \widehat{B A C}=\dfrac{B A}{A C}=\dfrac{3 a}{5 a}=\dfrac{3}{5}$.
			
		}{
			\begin{tikzpicture}[>=stealth,line join=round,line cap=round,font=\footnotesize,scale=1]
				\def\a{5};
				\def\b{2};
				\def\g{-40};
				\def\cao{4};
				\path
				(0,0) coordinate (A)
				(0:\a) coordinate (C)
				(\g:\b) coordinate (B)
				($(A)+(90:\cao)$) coordinate (S)
				($(S)!0.4!(B)$) coordinate (H)
				($(S)!0.6!(C)$) coordinate (K)
				;
				\draw (S)--(A)--(B)--(C)--(S)--(B) (A)--(H)--(K);
				\draw[dashed] (A)--(C) (A)--(K);
				\foreach \x/\g in {A/180,B/-90,C/0,S/90,H/45,K/30}\fill[black](\x) circle (1pt) +(\g:3mm) node {$\x$};
				\pic[draw,angle radius=2mm]{right angle=A--B--C};
				\pic[draw,angle radius=2mm]{right angle=A--H--S};
				\pic[draw,angle radius=2mm]{right angle=A--K--S};
			\end{tikzpicture}
		}
	}
\end{ex}

\begin{ex}%[1H3B4-3]%
	\immini[thm]{Hình chóp $S.ABC$ có $SA$ vuông góc với đáy và $SA=\dfrac{a}{2}$, tam giác $ABC$ vuông tại $A$, $AC=a\sqrt{3}$, $AB = a$. Tính góc giữa $(SBC)$ với $(ABC)$.
		\choice
		{$45^\circ$}
		{\True $30^{\circ}$}
		{$90^{\circ}$}
		{$60^{\circ}$}}{
		\begin{tikzpicture}[line cap=round,line join=round, >=stealth,scale=0.8]
			\def \xa{1.6}
			\def \xb{-1.5}
			\def \y{5}
			\def \z{3}
			\coordinate (A) at (0,0);
			\coordinate (B) at ($(A)+(\xa,\xb)$);
			\coordinate (C) at ($(A)+(\y,0)$);
			\coordinate (S) at ($ (A)+(0,\z) $);
			%	\coordinate (M) at ($(B)!0.35!(C)$);
			\draw [dashed] (A)--(C);
			\draw (S)--(A)--(B)--(C)--(S)--(B);
			\tkzDrawPoints(S,A,B,C)
			\tkzLabelPoints[right](C)
			\tkzLabelPoints[above](S)
			\tkzLabelPoints[left](A)
			\tkzLabelPoints[below](B)
			\tkzMarkRightAngles(C,A,B);
			\tkzMarkRightAngles(S,A,C);
	\end{tikzpicture}}
	\loigiai{
		\immini{
			Gọi $M$ là hình chiếu vuông góc của $A$ trên $BC$.\\
			$\heva{&BC \perp AM\\&BC \perp SA} \Rightarrow BC \perp (SAM)$.\\
			$\Rightarrow BM \perp SM$ (vì $SM \subset (SAM)$).\\
			Vì $\heva{&(SBC) \cap (ABC) =BC\\&SM \perp BC, SM \subset (SBC)\\&AM \perp BC, AM \subset (ABC)}$
		}{
			\begin{tikzpicture}[line cap=round,line join=round, >=stealth,scale=0.8]
				\def \xa{1.6}
				\def \xb{-1.5}
				\def \y{5}
				\def \z{3}
				\coordinate (A) at (0,0);
				\coordinate (B) at ($(A)+(\xa,\xb)$);
				\coordinate (C) at ($(A)+(\y,0)$);
				\coordinate (S) at ($ (A)+(0,\z) $);
				\coordinate (M) at ($(B)!0.35!(C)$);
				\draw [dashed] (M)--(A)--(C);
				\draw (M)--(S)--(A)--(B)--(C)--(S)--(B);
				\tkzDrawPoints(S,A,B,C,M)
				\tkzLabelPoints[right](C)
				\tkzLabelPoints[above](S)
				\tkzLabelPoints[left](A)
				\tkzLabelPoints[below](B,M)
				\tkzMarkRightAngles(C,A,B);
				\tkzMarkRightAngles(S,A,C);
				\tkzMarkAngle[size = 0.7](S,M,A)
		\end{tikzpicture}}
		$\Rightarrow ((SBC),(ABC))=(SM,AM) = \widehat{SMA}$ (vì tam giác $SAM$ vuông tại $A$).\\
		Tam giác $ABC$ vuông tại $A$ suy ra $\dfrac{1}{AM^2}=\dfrac{1}{AB^2}+\dfrac{1}{AC^2}=\dfrac{1}{a^2}+\dfrac{1}{(a\sqrt{3})^2}$ $\Rightarrow AM=\dfrac{a\sqrt{3}}{2}$.\\
		$\tan \widehat {SMA}=\dfrac{SA}{AM}=\dfrac{\dfrac{a}{2}}{\dfrac{a\sqrt{3}}{2}}=\dfrac{1}{\sqrt{3}} \Rightarrow \widehat {SMA}=30^{\circ}$.\\
		Vậy $((SBC),(ABC))=30^{\circ}$.
	}
\end{ex}

\begin{ex}%[1H3B4-3]%[Nguyễn Văn Nay]%
	Cho tứ diện $OABC$ có $OA,OB,OC$ đôi một vuông góc và $OB=OC=a\sqrt{6}$, $OA=a$ . Tính góc giữa hai mặt phẳng $(ABC)$ và $(OBC)$.
	\choice
	{$45^{\circ}$}
	{$90^{\circ}$}
	{$60^{\circ}$}
	{\True $30^{\circ}$}
	\loigiai
	{\immini{Trong $\triangle OBC$ kẻ đường cao $OH$.\\
			Theo đề bài ta có $\heva{&OA\perp OB\\&OA\perp OC} \Rightarrow OA \perp BC$.\\
			Lại có $BC \perp OH$ nên $BC \perp AH$.\\
			Do đó góc giữa $(ABC)$ và $(OBC)$ là góc $\widehat{AHO}$.\\
			Trong tam giác $OBC$ vuông tại $O$ có $OH$ là đường cao nên $$\dfrac{1}{OH^2}=\dfrac{1}{OB^2}+\dfrac{1}{OC^2}=\dfrac{1}{6a^2}+\dfrac{1}{6a^2}=\dfrac{1}{3a^2} \Rightarrow OH=a\sqrt{3}.$$
			Từ đó $\tan\widehat{AHO}=\dfrac{OA}{OH}=\dfrac{a}{a\sqrt{3}}=\dfrac{\sqrt{3}}{3}\Rightarrow\widehat{AHO}=30^{\circ}$.}{
			\begin{tikzpicture}[scale=0.6, font=\footnotesize, line join=round, line cap=round, >=stealth]
				\tikzset{label style/.style={font=\footnotesize}}
				\tkzDefPoints{0/0/O,7/0/B,3/-3/C}
				\coordinate (A) at ($(O)+(0,5)$);
				\coordinate (H) at ($(B)!0.5!(C)$);
				\tkzDrawPolygon(A,O,C,B)
				\tkzDrawSegments(A,C A,H)
				\tkzDrawSegments[dashed](O,B O,H)
				\tkzDrawPoints(A,O,B,C,H)
				\tkzLabelPoints[above](A)
				\tkzLabelPoints[below](C)
				\tkzLabelPoints[left](O)
				\tkzLabelPoints[right](B,H)
				%\tkzLabelSegment[below](O,C){$a\sqrt{6}$}
				%\tkzLabelSegment[above](O,B){$a\sqrt{6}$}
			\end{tikzpicture}
		}
	}
\end{ex}


\begin{ex}%[1H3B4-3]
	\immini[thm]{Cho hình chóp đều $S.ABC$ có $AB=a$, $SA=2a$. Gọi $\varphi$ là góc giữa mặt bên và mặt đáy. Tính  $\cos\varphi$.
		\choice
		{\True $\dfrac{\sqrt{5}}{15}$}
		{$\dfrac{\sqrt{3}}{6}$}
		{$\dfrac{\sqrt{5}}{5}$}
		{$\dfrac{\sqrt{3}}{3}$}}{
		\begin{tikzpicture}[scale=0.7, font=\footnotesize, line join=round, line cap=round, >=stealth]
			\tkzDefPoints{0/0/A,1.2/-1.6/B,4.5/0/C}
			\tkzCentroid(A,B,C)\tkzGetPoint{G}
			\coordinate (S) at ($(G)+(0,3)$);
			\coordinate (M) at ($(B)!1/2!(C)$);
			\tkzDrawPolygon(A,B,C,S)
			\tkzDrawSegments(S,B S,M)
			\tkzDrawSegments[dashed](M,A A,C S,G)
			\tkzDrawPoints[fill=black](A,B,C,S,G,M)
			\tkzLabelPoints[above](S)
			\tkzLabelPoints[below](B,G,M)
			\tkzLabelPoints[left](A)
			\tkzLabelPoints[right](C)
	\end{tikzpicture}}
	\loigiai{
		\immini{Gọi $M$là trung điểm của $BC$, $G$ là trọng tâm tam giác $ABC$ thì $SG \perp (ABC)$.\\
			Ta có $\heva{&BC\perp SG\\&BC\perp AM}$ nên góc giữa $(SBC)$ và $(ABC)$ là góc giữa $SM$ và $AM$, hay là góc $\widehat{SMA}$.\\
			Tam giác $SBM$ vuông tại $M$ có $$SM=\sqrt{SB^2-BM^2}=\sqrt{4a^2-\dfrac{a^2}{4}}=\dfrac{a\sqrt{15}}{2},$$
			$$GM=\dfrac{1}{3}\cdot AM=\dfrac{1}{3}\cdot \dfrac{a\sqrt{3}}{2}=\dfrac{a\sqrt{3}}{6}.$$
		}{
			\begin{tikzpicture}[scale=1, font=\footnotesize, line join=round, line cap=round, >=stealth]
				\tkzDefPoints{0/0/A,1.2/-1.6/B,4.5/0/C}
				\tkzCentroid(A,B,C)\tkzGetPoint{G}
				\coordinate (S) at ($(G)+(0,3.5)$);
				\coordinate (M) at ($(B)!1/2!(C)$);
				\tkzDrawPolygon(A,B,C,S)
				\tkzDrawSegments(S,B S,M)
				\tkzDrawSegments[dashed](M,A A,C S,G)
				\tkzDrawPoints[fill=black](A,B,C,S,G,M)
				\tkzLabelPoints[above](S)
				\tkzLabelPoints[below](B,G,M)
				\tkzLabelPoints[left](A)
				\tkzLabelPoints[right](C)
			\end{tikzpicture}
		}
		\noindent Tam giác $SGM$ vuông tại $G$, do đó $\cos\widehat{SMG}=\dfrac{GM}{SM}=\dfrac{\dfrac{a\sqrt{3}}{6}}{\dfrac{a\sqrt{15}}{2}}=\dfrac{\sqrt{5}}{15}$.
	}
\end{ex}

\begin{ex}%[1H3B4-3]
	Cho hình chóp $S.ABCD$ có đáy $ABCD$ là hình vuông cạnh $a.$ Cạnh bên $SA=x$ và vuông góc với mặt phẳng $\left(ABCD\right).$ Xác định $x$ để hai mặt phẳng $\left(SBC\right)$ và $\left(SCD\right)$ tạo với nhau một góc $60^\circ.$
	\choice
	{$x=\dfrac{3a}{2}$}
	{\True $x=a$}
	{$x=2a$}
	{$x=\dfrac{a}{2}$}
	\loigiai{
		\immini{
			Từ $A$ kẻ $AH$ vuông góc với $SB\left(H\in SB\right).$\\
			Ta có $\heva{& SA\perp BC \\ & AB\perp BC \\}\Rightarrow BC\perp \left(SAB\right)\Rightarrow BC\perp AH$ mà $AH\perp SB$ suy ra $AH\perp \left(SBC\right).$\\
			Từ $A$ kẻ $AK$ vuông góc với $SD\left(K\in SD\right),$ tương tự, chứng minh được $AK\perp \left(SCD\right).$\\
			Khi đó $SC\perp \left(AHK\right)$ suy ra $\left(\left(SBC\right);\left(SCD\right)\right)=\left(AH;AK\right)=\widehat{HAK}=60^\circ.$ \\
			Lại có $\triangle SAB=\triangle SAD\Rightarrow AH=AK$ mà $\widehat{HAK}=60^0$ suy ra tam giác $AHK$ đều.
		}{
			\begin{tikzpicture}[scale=0.6, font=\footnotesize, line join=round, line cap=round,>=stealth]
				\tkzDefPoints{0/0/B, 6/0/C, 8.5/2.5/D}
				\coordinate (A) at ($(B)+(D)-(C)$);
				\coordinate (S) at ($(A)+(0,5)$);
				\coordinate (H) at ($(S)!.4!(B)$);
				\coordinate (K) at ($(S)!.4!(D)$);
				\tkzDrawPoints[fill=black](A,B,C,D,S,H,K)
				\tkzDrawSegments[dashed](S,A A,C)
				\tkzDrawPolygon(S,B,C,D)
				\tkzDrawPolygon[dashed](A,B,D)
				\tkzDrawPolygon[dashed](A,H,K)
				\tkzDrawSegments(S,C)
				\tkzLabelPoints[above left](H)
				\tkzLabelPoints[above](S)
				\tkzLabelPoints[below](B,C,A)
				\tkzLabelPoints[above right](D,K)
			\end{tikzpicture}
		}
		\noindent Tam giác $SAB$ vuông tại $S,$ có $\dfrac{1}{AH^2}=\dfrac{1}{SA^2}+\dfrac{1}{AB^2}\Rightarrow AH=\dfrac{xa}{\sqrt{x^2+a^2}}.$\\
		Suy ra $SH=\sqrt{SA^2-AH^2}=\dfrac{x^2}{\sqrt{x^2+a^2}}\Rightarrow \dfrac{SH}{SB}=\dfrac{x^2}{x^2+a^2}.$\\
		Vì $HK\parallel BD$ suy ra $\dfrac{SH}{SB}=\dfrac{HK}{BD}\Leftrightarrow \dfrac{x^2}{x^2+a^2}=\dfrac{xa}{\sqrt{x^2+a^2}\cdot a\sqrt{2}}\Leftrightarrow \dfrac{x}{\sqrt{x^2+a^2}}=\dfrac{1}{\sqrt{2}}\Rightarrow x=a.$
	}
\end{ex}
\Closesolutionfile{ans}

\ind{PHẦN II.} \inden{Câu trắc nghiệm đúng sai. Trong mỗi ý a), b), c), d) ở mỗi câu, học sinh chọn đúng hoặc sai.}\\
\Opensolutionfile{ans}[ans/1H7-B4-2]
\begin{ex}%[1H8H4-3]
	Cho hình chóp $S.ABC$ có $SA \perp \left( ABC \right)$ và đáy $ABC$ là tam giác đều. Xét tính đúng sai của các mệnh đề sau.
	\choiceTF
	{\True $\left( SAB \right)\perp \left( ABC \right)$}
	{\True $\left( SAC \right)\perp \left( ABC \right)$}
	{\True Gọi $H$ là trung điểm của cạnh $BC$. Khi đó $\widehat{AHS}$ là góc giữa hai mặt phẳng $\left( SBC \right)$ và $\left( ABC \right)$}
	{Góc giữa hai mặt phẳng $\left( SBC \right)$ và $\left( SAC \right)$ là $\widehat{ACB}$}
	\loigiai{
		\begin{center}
			\begin{tikzpicture}[line join=round,line cap=round,line width=.6pt,font=\footnotesize,scale=1]
				\coordinate[label=left:$A$] (A) at (0,0);
				\coordinate[label=below left:$B$] (B) at (2,-2);
				\coordinate[label=right:$C$] (C) at (4,0);
				\coordinate[label=above left:$S$] (S) at ($(A)+(90:3)$);
				
				\coordinate[label=below right:$H$] (H) at ($(B)!0.5!(C)$);
				\coordinate[label=above left:$K$] (K) at ($(A)!0.5!(C)$);
				
				\coordinate[label=above right:$F$] (F) at ($(C)!(K)!(S)$);
				\draw (A)--(B)--(C)--(S)--cycle (S)--(B) (S)--(H) (B)--(F);
				\draw[dashed] (A)--(C) (A)--(H) (B)--(K)--(F);
				\draw pic[draw, angle radius=3mm, angle eccentricity=1.5]{right angle = C--A--S};
				\draw pic[draw, angle radius=3mm, angle eccentricity=1.5]{right angle = B--A--S};
				\fill (A)circle(2pt) (B)circle(2pt) (C)circle(2pt) (S)circle(2pt) (H)circle(2pt) (K)circle(2pt) (F)circle(2pt);
				
			\end{tikzpicture}
			
		\end{center}
		\begin{itemchoice}
			\itemch Đúng. Vì $SA \perp \left( ABC \right)$ và $SA\subset \left( SAB \right)$ nên $\left( SAB \right) \perp \left( ABC \right)$.
			\itemch Đúng. Vì $SA \perp \left( ABC \right)$ và $SA\subset \left( SAC \right)$ nên $\left( SAC \right) \perp \left( ABC \right)$.
			\itemch Đúng. Vì $H$ là trung điểm của cạnh $BC$ nên $AH \perp BC$.\\
			Mặt khác $SA \perp \left( ABC \right), \,BC\subset \left( ABC \right)$ nên $SA \perp BC$.\\
			Do đó $BC \perp \left( SAH \right)\Rightarrow BC \perp SH$.\\
			Ta có $\left( SBC \right)\cap \left( ABC \right)=BC$\\
			$AH\perp BC,\,AH\subset \left( ABC \right)$.\\
			$SH\perp BC,\,SH\subset \left( ABC \right)$.\\
			Suy ra góc giữa hai mặt phẳng $\left( SBC \right)$ và $\left( ABC \right)$ là góc giữa hai đường thẳng $AH$ và $SH$ hay góc $\widehat{AHS}$.
			\itemch Sai. Gọi $K$ là trung điểm của cạnh $AC$ $\Rightarrow BK \perp AC$\\
			Mà $SA \perp \left( ABC \right),\,BK\subset \left( ABC \right)$ nên $SA \perp BK$.\\
			Do đó $BK\perp \left( SAC \right)\Rightarrow BK \perp SC\,\,\left( 1 \right)$.\\
			Trong mặt phẳng $\left( SAC \right)$ kẻ $KF \perp SC\,\,\left( F\in SC \right)\,\,\left( 2 \right)$.\\
			Từ (1) và (2) suy ra $SC\perp \left( BKF \right)\Rightarrow BF\perp SC$.\\
			Ta có: $\left( SAC \right)\cap \left( SBC \right)=SC$\\
			$KF\perp SC,\,KF\subset \left( SAC \right)$\\
			$BF\perp SC,\,BF\subset \left( SBC \right)$\\
			Suy ra góc giữa hai mặt phẳng $\left( SBC \right)$ và $\left( SAC \right)$ là góc giữa hai đường thẳng $KF$ và $BF$, hay góc $\widehat{BFK}$.
		\end{itemchoice}
	}
\end{ex}
%Câu 2:
\begin{ex}%[1H8H4-3]
	Cho lăng trụ $ABC.A'B'C'$ có đáy là tam giác đều cạnh $a$, $A'A=A'B=A'C=a$. $O$ là tâm đường tròn ngoại tiếp tam giác $ABC$, $H$ là trung điểm $BC$. Xét tính đúng sai của các mệnh đề sau.
	\choiceTF
	{\True Góc giữa hai mặt phẳng $\left( A'AH \right)$ và $\left( ABC \right)$ bằng $90^\circ $}
	{\True Góc giữa hai mặt phẳng $\left( A'AH \right)$ và $\left( BCC'B' \right)$ bằng $90^\circ $}
	{Góc giữa hai mặt phẳng $\left( A'BC \right)$ và $\left( ABC \right)$ bằng $90^\circ $}
	{\True Gọi $\alpha $ là góc giữa mặt phẳng $\left( A'BC \right)$ và $\left( ABC \right)$. Khi đó $\tan \alpha =2\sqrt 2$}
	\loigiai{ 
		\begin{center} 
			\begin{tikzpicture}[line join=round,line cap=round,line width=.6pt,font=\footnotesize,scale=1]
				\coordinate[label=left:$A$] (A) at (0,0);
				\coordinate[label=below left:$B$] (B) at (1,-1);
				\coordinate[label=right:$C$] (C) at (4,0);
				\coordinate[label=below right:$H$] (H) at ($(B)!.5!(C)$);
				\coordinate[label=below:$O$] (O) at ($(A)!2/3!(H)$);
				\coordinate[label=left:$A'$] (A1) at ($(O)+(90:4)$);
				\coordinate[label=above :$B'$] (B1) at ($(B)-(A)+(A1)$);
				\coordinate[label=right:$C'$] (C1) at ($(C)-(A)+(A1)$);
				\draw (A1)--(A)--(B)--(C)--(C1)--(A1)--(B1)--(C1) (B)--(B1);
				\draw[dashed] (A)--(C) (O)--(A1) (A)--(H)--(A1)--(C);
				\draw ($ (O)!5pt!(C)$)--($(O)!2!($($(O)!5pt!(C)$)!.5!($(O)!5pt!(A1)$)$)$)--($(O)!5pt!(A1)$);
				\fill (A)circle(2pt) (B)circle(2pt) (C)circle(2pt) (A1)circle(2pt) (B1)circle(2pt) (C1)circle(2pt) (H)circle(2pt) (O)circle(2pt);
			\end{tikzpicture} 
		\end{center}
		\begin{itemchoice}
			\itemch Đúng. Ta có $A'A=A'B=A'C$ và $O$ là tâm đường tròn ngoại tiếp tam giác $ABC$ nên $A'O\perp \left( ABC \right)$. Do $A'O\subset \left( A'AH \right)$ nên $\left( A'AH \right)\perp \left( ABC \right)$.
			\itemch Đúng. Ta có $\heva{
				& A'H\perp BC \ 
				& AH\perp BC \ 
			}\Rightarrow BC\perp \left( A'AH \right)$.\\
			Do $BC\subset \left( BCC'B' \right)$ nên $\left( A'AH \right)\perp \left( BCC'B' \right)$.
			\itemch Sai. Ta có $BC\perp \left( A'AH \right)\Rightarrow BC\perp A'H$, tam giác $ABC$ đều nên $BC\perp AH$.\\
			Khi đó góc giữa $\left( A'BC \right)$ và $\left( ABC \right)$ là góc giữa $A'H$ và $AH$. Do tam giác $A'OH$ vuông tại $O$ nên $\widehat{A'HA}\ne 90^\circ $, vậy $\left( A'H,AH \right)\ne 90^\circ $.
			\itemch Đúng. Ta có tam giác $ABC$ đều cạnh $a$ nên $AH=\dfrac{a\sqrt 3} 2$, $AO=\dfrac{a\sqrt 3} 3$, $OH=\dfrac{a\sqrt 3} 6$.\\
			Tam giác $A'OA$ vuông tại $O$ nên $A'O=\sqrt{A'a^2-AO^2}=\sqrt{a^2-\dfrac{a^2}{3}}=\dfrac{a\sqrt 6} 3$.\\
			Góc giữa $\left( A'BC \right)$ và $\left( ABC \right)$ là góc $\alpha =\widehat{A'HO}$. Khi đó $\tan \alpha =\dfrac{A'O}{OH}=\dfrac{\dfrac{a\sqrt 6} 3}{\dfrac{a\sqrt 3} 6}=2\sqrt 2$.
		\end{itemchoice}
	}
\end{ex}
%Câu 3:
\begin{ex}%[1H8H4-3]
	Cho hình chóp $SABCD$ có đáy $ABCD$ là hình vuông tâm $O$ cạnh $2a$, $SO\perp \left( ABCD \right)$ và $SO=a$. Gọi $I$, $J$ là trung điểm của $AD$, $BC$. Xét tính đúng sai của các mệnh đề sau.
	\choiceTF
	{\True $\left( SAC \right)\perp \left( ABCD \right)$}
	{\True ${(SAC) }$ $\perp \left( SBD \right)$}
	{\True $\left( SIJ \right)\perp \left( SBC \right)$}
	{$\bigl( \left( SAD \right),\left( SBC \right) \bigr)=60^\circ $}
	\loigiai{
		\begin{center}
			\begin{tikzpicture}[line join=round,line cap=round,line width=.6pt,font=\footnotesize,scale=1]
				\coordinate[label=below left:$B$] (B) at (0,0);
				\coordinate[label=above right:$A$] (A) at (1,.8);
				\coordinate[label=below right:$C$] (C) at (4,0);
				\coordinate[label=above right:$D$] (D) at ($(C)-(B)+(A)$);
				\coordinate[label=below:$O$] (O) at ($(A)!.5!(C)$);
				\coordinate[label=above left:$S$] (S) at ($(O)+(90:4)$);
				\coordinate[label=above left:$I$] (I) at ($(A)!0.5!(D)$);
				
				\coordinate[label=below left:$J$] (J) at ($(B)!0.5!(C)$);
				
				\draw (B)--(C)--(D)--(S)--cycle (J)--(S)--(C);
				\draw[dashed] (C)--(A)--(D)--(B) (O)--(S)--(A)--(B) (S)--(I)--(J);
				\fill (A)circle(2pt) (B)circle(2pt) (C)circle(2pt) (D)circle(2pt) (S)circle(2pt) (O)circle(2pt) (I)circle(2pt) (J)circle(2pt);
			\end{tikzpicture}
			
		\end{center}
		\begin{itemchoice}
			\itemch Đúng. Do $SO\perp \left( ABCD \right)\Rightarrow \left( SAC \right)\perp \left( ABCD \right)$.
			\itemch Đúng. Ta có: $ABCD$ là hình vuông suy ra $AC\perp BD~\left( 1 \right)$\\
			Hơn nữa $SO\perp \left( ABCD \right)\Rightarrow SO\perp AC\left( 2 \right)$\\
			Từ (1) và (2) suy ra $AC\perp \left( SBD \right)\Rightarrow \left( SAC \right)\perp \left( SBD \right)$.
			\itemch Đúng. Do $SO\perp \left( ABCD \right)\Rightarrow SO\perp BC~\left( 3 \right)$\\
			Vì $I$, $J$ là trung điểm của $AD,BC\Rightarrow IJ\perp BC\left( 4 \right)$\\
			Từ (3) và (4) suy ra $BC\perp \left( SIJ \right)\Rightarrow \left( SBC \right)\perp \left( SIJ \right)$.
			\itemch Sai. Trong tam giác $SIJ$ ta thấy $SO=\dfrac{1}{2}IJ\Rightarrow \Delta SIJ$ vuông tại $S\Rightarrow SI\perp SJ\left( 5 \right)$\\
			Mà theo chứng minh trên ta có $BC\perp \left( SIJ \right)\Rightarrow BC\perp SI\left( 6 \right)$\\
			Từ (5) và (6) suy ra $SI\perp \left( SBC \right)\Rightarrow \left( SAD \right)\perp \left( SBC \right)\Rightarrow \bigl( \left( SAD \right),\left( SBC \right) \bigr)=90^\circ.$
		\end{itemchoice}
	}
\end{ex}
%Câu 4:
\begin{ex}%[1H8H4-3]
	Cho hình chóp $S.ABCD$ có đáy $ABCD$ là hình thang vuông tại $A$ có $AB=2a$, $ AD=DC=a$, $SA=a$ và $SA\perp \left( ABCD \right)$. Xét tính đúng sai của các mệnh đề sau.
	\choiceTF
	{\True $\left( SAC \right)\perp \left( ABCD \right)$}
	{\True $\left( SAC \right)\perp \left( SBC \right)$}
	{$CB\perp \left( SAD \right)$}
	{\True $\sin \bigl( \left( SAB \right),\left( SBC \right) \bigr)=\dfrac 5{\sqrt{30}}$}
	\loigiai{
		\begin{center}
			\begin{tikzpicture}[scale=0.8, font=\footnotesize,line join=round, line cap=round, >=stealth]
				\coordinate (A) at (0,0);
				\coordinate (D) at (2,-2);
				\coordinate (B) at (10,0);
				\coordinate (C) at ($(D)+(5,0)$);
				\coordinate (I) at ($(A)!0.5!(B)$);
				\coordinate (S) at ($(A)+(0,4)$);
				\coordinate (M) at ($(B)!(I)!(S)$);
				\foreach \i in {D,C,B}{\draw (S)--(\i);}
				\draw (B)--(C)--(D) (C)--(M) (S)--(A)--(D);
				\draw[dashed,thin](B)--(A)--(D) (C)--(I)--(M) (A)--(C) (I)--(D);
				\coordinate (O) at (intersection of A--C and D--I); 
				\foreach \i/\g in {S/90,A/-135,D/-90,C/-90,B/-90,I/120,M/40,O/90}{\draw[fill=black](\i) circle (1.5pt) ($(\i)+(\g:3mm)$) node[scale=1]{$\i$};}
			\end{tikzpicture} 
		\end{center}
		\begin{itemchoice}
			\itemch Đúng. Vì $SA\perp \left( ABCD \right)\Rightarrow \left( SAC \right)\perp \left( ABCD \right).$
			\itemch Đúng. Gọi $I$ là trung điểm của $AB$. Ta có $AICD$ là hình vuông nên $AC\perp DI$ (1).\\
			Mặt khác tứ giác $IBCD$ là hình bình hành nên $ID\parallel BC$ (2).\\
			Từ (1) và (2) ta có $AC\perp BC$ (3).\\
			Do $SA\perp \left( ABCD \right)\Rightarrow SA\perp CB$ (4).\\
			Từ (3) và (4) ta có $CB\perp \left( SAC \right)\Rightarrow \left( SBC \right)\perp \left( SAC \right).$
			\itemch Sai. Do $CD\perp \left( SAD \right)$ nên $CB$ không thể vuông góc với $\left( SAD \right)$
			\itemch Đúng. Ta có $I$ là trung điểm của $AB\Rightarrow CI \parallel AD\Rightarrow CI\perp AB$ (5).\\
			Do $SA\perp \left( ABCD \right)\Rightarrow SA\perp CI$ (6)\\
			Từ (5) và (6) suy ra $CI\perp \left( SAB \right)\Rightarrow CI\perp SB$. Từ $I$ hạ $IM\perp SB,M\in SB.$\\
			Ta có $IM\perp SB$ và $CI\perp SB$, nên $SB\perp \left( ICM \right)\Rightarrow SB\perp CM.$\\
			Do $\heva{&\left( SAB \right)\cap \left( SBC \right)=SB \\&
				IM\perp SB \\&
				CI\perp SB} \Rightarrow \bigl( \left( SAB \right),\left( SBC \right) \bigr)=\left( IM,CM \right)=\widehat{IMC}$.\\
			Theo phần a, ta có $CB\perp \left( SAC \right)\Rightarrow CB\perp SC$.\\
			Vậy tam giác $SCB$ vuông tại $C$.\\
			Ta có $SC=a\sqrt 3$, $CB=a\sqrt 2$.\\
			Xét tam giác $SCB$ vuông tại $C$ ta có
			$$\dfrac 1{CM^2}=\dfrac 1{SC^2}+\dfrac 1{Cb^2}\Rightarrow CM=\dfrac{SC.CB}{SC+CB}=\dfrac{a\sqrt{30}} 5.$$
			Xét tam giác $MIC$ vuông tại $I$, ta có
			$$\sin\widehat{IMC}=\dfrac{CI}{CM}=\dfrac{5a}{a\sqrt{30}}=\dfrac 5{\sqrt{30}}\Rightarrow \sin \bigl( \left( SAB \right),\left( SBC \right) \bigr)=\dfrac 5{\sqrt{30}}.$$
		\end{itemchoice}
	}
\end{ex}
\Closesolutionfile{ans}

\ind{PHẦN III.} \inden{Câu trắc nghiệm trả lời ngắn.}\\
\Opensolutionfile{ans}[ans/1H7-B4-3]

\begin{ex}
	Cho hình chóp ${S.ABCD}$ có đáy là hình chữ nhật, $SC\bot (ABCD)$, $CD=a, SB=\sqrt{21} a, SC=\sqrt{5} a$. Tính số đo góc phẳng nhị diện ${[C,AB,S]}$ (tính theo đơn vị độ, làm tròn đến hàng phần chục).\\
	\shortans[oly]{${29,2}$}
	\loigiai{ 
		Ta có: $CB \bot AB, AB \bot SC \Rightarrow AB \bot SB$. Do đó ${[C,AB,S]} =\widehat{SBC}$.
		
		$\displaystyle \sin \widehat{SBC}=\dfrac{SC} {SB}= \dfrac{\sqrt{5} a} {\sqrt{21} a}=\frac{\sqrt{105}}{21}$.
		Suy ra $\widehat{SBC} = 29,2^\circ$. 
}\end{ex}

\begin{ex}
	Cho hình chóp ${S.ABCD}$ có đáy là hình chữ nhật, $SD\bot (ABCD)$, $DA=5 a, DC=9 a, SD=6 a$. Tính số đo góc phẳng nhị diện ${[D,AC,S]}$ (tính theo đơn vị độ, làm tròn đến hàng phần chục).\\
	\shortans[oly]{${53,9}$}
	\loigiai{ 
		Dựng $DH \bot AC,AC \bot SD \Rightarrow SH \bot AC$. Do đó ${[D,AC,S]} =\widehat{SHD}$.
		
		$\displaystyle AH= \dfrac{DA.DC} {\sqrt{DA^2 + DC^2}}= \frac{45 \sqrt{106}}{106}a$.
		
		$\displaystyle \tan \widehat{SHD}=\dfrac{SD} {DH}= \dfrac{6 a} {\frac{45 \sqrt{106} a}{106}}=\frac{2 \sqrt{106}}{15}$.
		Suy ra $\widehat{SHD} = 53,93^\circ$. 
}\end{ex}

\begin{ex}%[1H3B4-3]
	\immini[thm]{Cho hình chóp $S.ABCD$ có đáy $ABCD$ là hình vuông, $SA \perp (ABCD)$, $SA=AB$. Góc giữa $(SBC)$ và $(ABCD)$ có số đo bằng bao nhiêu? (tính theo đơn vị độ)\\
	\shortans[oly]{${45}$}
}{
		\begin{tikzpicture}[line join=round,line cap=round,line width=.6pt,font=\footnotesize,scale=0.7]
			\coordinate[label=below left:$B$] (B) at (0,0);
			\coordinate[label=above left:$A$] (A) at (1,.8);
			\coordinate[label=below right:$C$] (C) at (4,0);
			\coordinate[label=above right:$D$] (D) at ($(C)-(B)+(A)$);
			\coordinate[label=above left:$S$] (S) at ($(A)+(90:3)$);
			\draw (B)--(C)--(D)--(S)--cycle (S)--(C);
			\draw[dashed] (A)--(D) (S)--(A)--(B);
			\draw ($ (A)!5pt!(D)$)--($(A)!2!($($(A)!5pt!(D)$)!.5!($(A)!5pt!(S)$)$)$)--($(A)!5pt!(S)$);
			\fill (A)circle(1.5pt) (B)circle(1.5pt) (C)circle(1.5pt) (D)circle(1.5pt) (S)circle(1.5pt);
	\end{tikzpicture}}
	\loigiai{
		\immini{Ta có $BC\perp (SAB) \Rightarrow BC \perp SB.\hfill(1)$\\
			Mặt phẳng $(SBC)$ và $(ABCD)$ có giao tuyến là $BC$. Vì $ABCD$ là hình vuông nên $BC \perp AB. \hfill (2)$\\
			Từ $(1)$ và $(2)$ suy ra góc giữa $(SBC)$ và $(ABCD)$ là góc $\widehat{(SB,BA)}=\widehat{SBA}=45^\circ$ (vì $\triangle SAB$ vuông cân tại $A$).}
		{\begin{tikzpicture}[line join=round,line cap=round,line width=.6pt,font=\footnotesize,scale=0.7]
				\coordinate[label=below left:$B$] (B) at (0,0);
				\coordinate[label=above left:$A$] (A) at (1,.8);
				\coordinate[label=below right:$C$] (C) at (4,0);
				\coordinate[label=above right:$D$] (D) at ($(C)-(B)+(A)$);
				\coordinate[label=above left:$S$] (S) at ($(A)+(90:4)$);
				\draw (B)--(C)--(D)--(S)--cycle (S)--(C);
				\draw[dashed] (A)--(D) (S)--(A)--(B);
				\draw ($ (A)!5pt!(D)$)--($(A)!2!($($(A)!5pt!(D)$)!.5!($(A)!5pt!(S)$)$)$)--($(A)!5pt!(S)$);
				\fill (A)circle(1.5pt) (B)circle(1.5pt) (C)circle(1.5pt) (D)circle(1.5pt) (S)circle(1.5pt);
			\end{tikzpicture}
		}
	}
\end{ex}

\begin{ex}%[1H3B4-3]
	\immini[thm]{Cho hình chóp đều $S.ABCD$ có tất cả các cạnh đều bằng $a$. Gọi $\varphi $ là góc giữa hai mặt phẳng $\left(SBD\right)$ và $\left(SCD\right)$. Tính $\tan \varphi$ (làm tròn đến hàng phần chục).\\
	\shortans[oly]{${1,4}$}
	}{
		\begin{tikzpicture}[scale=1, line join=round, line cap=round]
			\tkzDefPoints{0/0/A,-2/-1.6/B,1.6/-1.6/C}
			\coordinate (D) at ($(A)+(C)-(B)$);
			\coordinate (O) at ($(A)!1/2!(C)$);
			\coordinate (S) at ($(O)+(0,3)$);
			\tkzDrawPolygon(S,B,C,D)
			\tkzDrawSegments(S,C)
			\tkzDrawSegments[dashed](A,S A,B A,D A,C B,D S,O)
			\tkzDrawPoints[fill=black,size=4](D,C,A,B,S)
			\tkzLabelPoints[above](S)
			\tkzLabelPoints[below](A,B,C,O)
			\tkzLabelPoints[right](D)
	\end{tikzpicture}}
	\loigiai{
		\immini{
			Gọi $O=AC\cap BD$. Do hình chóp $S.ABCD$ đều nên $SO\perp \left(ABCD\right)$.\\
			Gọi $M$ là trung điểm của $SD$. Tam giác $SCD$ đều nên $CM\perp SD$.\\
			Tam giác $SBD$ có $SB=SD=a$, $BD=a\sqrt{2}$ nên vuông tại $S\Rightarrow SB\perp SD\Rightarrow OM\perp SD.$\\
			Do đó $\left(\left(SBD\right),\left(SCD\right)\right)=(OM,CM)=\widehat{OMC}$.\\
			Ta có $\heva{& OC\perp BD \\ & OC\perp SO \\}\Rightarrow OC\perp \left(SBD\right)\Rightarrow OC\perp OM$.\\
			Tam giác vuông $MOC$, có $\tan \widehat{CMO}=\dfrac{OC}{OM}=\sqrt{2}.$
		}{
			\begin{tikzpicture}[scale=0.7, font=\footnotesize, line join=round, line cap=round,>=stealth]
				\tkzDefPoints{0/0/B, 6/0/C, 3/2/A}
				\coordinate (D) at ($(A)+(C)-(B)$);
				\coordinate (O) at ($(B)!.5!(D)$);
				\coordinate (S) at ($(O)+(0,6)$);
				\coordinate (M) at ($(S)!.5!(D)$);
				\tkzDrawPoints[fill=black](A,B,C,D,S,O,M)
				\tkzDrawSegments[dashed](S,A B,A A,D S,O A,C B,D O,M)
				\tkzDrawSegments(S,B S,D S,C B,C D,C C,M)
				\tkzLabelPoints[right](D)
				\tkzLabelPoints[above](S,M)
				\tkzLabelPoints[left](A)
				\tkzLabelPoints[below](B,C,O)
				\tkzMarkRightAngles(A,O,B S,O,D C,M,D O,M,S)
			\end{tikzpicture}
		}
	}
\end{ex}

\begin{ex}
	Trong hình chóp tam giác đều có góc giữa cạnh bên và mặt đáy bằng $60^{\circ}$.  Gọi $\varphi $ là góc giữa mặt phẳng bên với mặt đáy. Tính $\tan \varphi$ (làm tròn đến hàng phần chục).\\
	\shortans[oly]{${3,5}$}
	\loigiai
	{
		\immini
		{
			Kí hiệu hình chóp $S.ABC$ là hình chóp tam giác đều.\\
			Gọi $M$ là trung điểm của $BC$, $G$ là trọng tâm của tam giác $ABC$.\\
			Khi đó $SG \perp (ABC)$, do đó góc giữa $SC$ với $(ABC)$ bằng góc giữa $SC$ và $GC$ chính là $\widehat{SCG} = 60^\circ$.\\
			Đặt $AB=a$. Ta có $AM = \dfrac{a\sqrt{3}}{2}$, $GC = \dfrac{a\sqrt{3}}{3}$, $GM = \dfrac{a\sqrt{3}}{6}$.\\
			Trong tam giác vuông $SGC$ ta có $$\tan\widehat{SCG} = \dfrac{SG}{GC} \Rightarrow SG = GC\cdot \tan\widehat{SCG} = \dfrac{a\sqrt{3}}{3} \cdot \sqrt{3} = a.$$
			Ta có $BC\perp AM$ và $BC\perp SG$ nên $BC \perp (SAM)$, suy ra $BC \perp SM$.\\
			Lại có $(SBC) \cap (ABC) = BC$.
		}
		{
			\begin{tikzpicture}[line join=round, line cap=round, >=stealth,font=\footnotesize, scale=1]
				\tkzDefPoints{0/0/A, 4/0/C, 2.5/-1.5/B}
				\tkzCentroid(A,B,C)\tkzGetPoint{G}
				\tkzDefShiftPoint[G](90:4){S}
				\tkzDefMidPoint(B,C)\tkzGetPoint{M}
				\tkzDrawPoints[fill=black, size=3pt](A,B,C,G,S,M)
				\tkzDrawSegments(S,A S,B S,C S,M A,B B,C)
				\tkzDrawSegments[dashed](S,G A,C A,M G,C)
				\tkzMarkSegments[mark=|, size=0.1cm](B,M C,M)
				\tkzMarkRightAngles[size=0.15](S,G,A A,M,B S,M,C)
				\node[above] at (S){$S$};
				\node[below] at (G){$G$};
				\node[below] at (B){$B$};
				\node[left] at (A){$A$};
				\node[right] at (C){$C$};
				\node[right] at (M){$M$};
			\end{tikzpicture}
		}
		\noindent
		Vậy góc giữa $(SBC)$ và $(ABC)$ bằng góc giữa $SM$ và $AM$, chính là $\widehat{SMG}$.\\
		Trong tam giác vuông $SMG$ ta có $\tan \widehat{SMG} = \dfrac{SG}{GM} = \dfrac{a}{\dfrac{a\sqrt{3}}{6}} = \dfrac{6}{\sqrt{3}} = 2\sqrt{3} \approx 3,5$.
	}
\end{ex}

\begin{ex}%[Thi thử, Ngô Sĩ Liên - Bắc Giang, 2019]%[Trần Nhân Kiệt, 12EX6-2019]%[1H3B4-3]%
	\immini[thm]{Cho hình lập phương $ABCD.A'B'C'D'$ có cạnh bằng $a$. Gọi $O'$ là tâm của hình vuông $A'B'C'D'$ và $\alpha$ là góc giữa hai mặt phẳng $(O'AB)$ và $(ABCD)$. Tính $\tan \varphi$ (làm tròn đến hàng phần chục).\\
		\shortans[oly]{${2,0}$}
	}{
		\begin{tikzpicture}[scale=0.8, line join=round, line cap=round]
			\tkzDefPoints{0/0/A,-1.3/-1.1/B,2/-1.1/C}
			\coordinate (D) at ($(A)+(C)-(B)$);
			\coordinate (A') at ($(A)+(0,2.5)$);
			\tkzDefPointsBy[translation=from A to A'](B,C,D){B'}{C'}{D'}
			\tkzDrawPolygon(A',B',B,C,D,D')
			\tkzDrawSegments(B',C' C',D' C,C')
			\tkzDrawSegments[dashed](A,B A,D A,A')
			\tkzDrawPoints[fill=black,size=4](A,B,D,C,A',B',C',D')
			\tkzLabelPoints[above](A',D')
			\tkzLabelPoints[below](A,B,C)
			\tkzLabelPoints[left](B')
			\tkzLabelPoints[right](C',D)
	\end{tikzpicture}}
	\loigiai{
		\immini
		{
			Gọi $O$, $M$ lần lượt là trung điểm $AC$, $AB$.\\
			Ta có $AB\perp (OMO')$\\
			$\Rightarrow $ góc giữa $(O'AB)$ và $(ABCD)$ là $\alpha=\widehat{OMO'}$.\\
			Xét tam giác $OM=\dfrac{a}{2}$, $OO'=a$\\
			$\Rightarrow \tan\alpha=\dfrac{OO'}{OM}=\dfrac{a}{\tfrac{a}{2}}=2$.
		}
		{
			\begin{tikzpicture}[scale=1, font=\footnotesize, line join=round, line cap=round, >=stealth]
				\tkzDefPoints{0/0/A,-1.5/-1.5/B,2/-1.5/C}
				\coordinate (D) at ($(A)+(C)-(B)$);
				\coordinate (A') at ($(A)+(0,3)$);
				\tkzDefPointsBy[translation=from A to A'](B,C,D){B'}{C'}{D'}
				\coordinate (O') at ($(A')!0.5!(C')$);
				\coordinate (O) at ($(A)!0.5!(C)$);
				\coordinate (M) at ($(A)!0.5!(B)$);
				\tkzDrawPolygon(A',B',B,C,D,D')
				\tkzDrawSegments(B',C' C',D' C,C' A',C' B',D')
				\tkzDrawSegments[dashed](A,B A,D A,A' O',A O',B A,C B,D O,O' O,M O',M)
				\tkzDrawPoints[fill=black,size=4](A,B,D,C,A',B',C',D',O',O,M)
				\tkzLabelPoints[above](A',D',O')
				\tkzLabelPoints[below](A,B,C,O,M)
				\tkzLabelPoints[left](B')
				\tkzLabelPoints[right](C',D)
			\end{tikzpicture}
		}
	}
\end{ex}

\Closesolutionfile{ans}

